\subsection{معرفی و انگیزه}
این مکانیزم بر مبنای مبانی نظری پیشرفته در ادبیات توجه و فشرده‌سازی \cite{ref_zerosattention_2024} توسعه یافته است.

یکی از مشکلات بنیادین در مکانیزم‌های توجه خطی کِرنلی استاندارد (مانند کِرنل‌های \lr{ELU+1} یا \lr{ReLU})، وجود سوگیری غیرصفر در بازنمایی بردارها و انباشت پتانسیل مقدار ثابت در طول دنباله است که به «پدیده غلبه مؤلفه مشترک» معروف است. این پدیده سبب می‌شود که با افزایش طول دنباله، تمایزپذیری میان توکن‌ها در فضای توجه خطی به شدت کاهش یافته و خروجی به سمت بردار میانگین توکن‌های گذشته میل کند.

مکانیزم توجه \lr{Zero-Mean Linear Attention (ZeroS Attention)} با هدف بازتولید خاصیت تمایزپذیری توجه چگال در ساختار محاسباتی خطی معرفی شده است. ایده بنیادین در \lr{ZeroS}، صفرسازی میانگین بردارهای کلید (\lr{Mean-Centering}) پیش از اعمال نگاشت فعال‌سازی نامنفی \lr{ReLU} است. با کسر کردن میانگین هندسی کلیدها، مؤلفه همبستگی همگانی حذف شده و تنها الگوهای تفاضلی و متمایزکننده توکن‌ها در ماتریس‌های حافظه خطی ذخیره می‌شوند.
